Valuing American options along simulated stock paths requires an implied-volatility (IV) input at every date, which a stock-return model alone does not provide. In this study, we developed a simulator that assigned an IV to American-style option contracts on simulated dates for any chosen stock model. We fitted parametric and neural IV surfaces, which predict IV from moneyness and time to expiration, to vendor option quotes for 31 tickers. Along stock paths from a jump hidden Markov model with capped daily returns, every option's implied variance moved partway each day toward its surface prediction, following the Heston variance equation with random shocks that tended to raise IV when the stock fell, and a binomial tree converted that IV into option values and price sensitivities. Neural surfaces fitted by sector or ticker matched the quotes more closely than one parametric surface, although their errors varied by date and carried into dollar prices. Repricing identical stock paths under five ways of updating IV changed option values before expiration and the worst losses from closing short positions early, but not payoffs at expiration. In forecasts of later Goldman Sachs and Eli Lilly option prices, stochastic IV did little better than fixed IV, and rerunning the forecasts with the realized stock paths showed that stock-forecast errors caused most misses. The hidden Markov model and two simpler stock models tuned on earlier data all predicted stock prices about as accurately as assuming no change, and the one whose volatility adapted to recent returns improved predicted price ranges for Eli Lilly but not Goldman Sachs. The framework supported reproducible comparisons of IV and stock assumptions but did not forecast better than simpler alternatives, and incomplete option histories, overlapping forecast dates, and simulated prices inconsistent across strikes limited broader use.
